\documentclass{article}
\usepackage[utf8]{inputenc}
\usepackage{geometry}
\geometry{a4paper, margin=1in}
\usepackage{hyperref}
\usepackage{titlesec}
\usepackage{enumitem}
\usepackage{xcolor}

\title{Codebase Evolution Report: EvoBrain-main vs EvoBrain}
\author{Antigravity Analysis}
\date{\today}

\begin{document}
\maketitle

\section*{Executive Summary}
The updated \texttt{EvoBrain} repository introduces significant improvements over the original \texttt{EvoBrain-main}. The primary focus of the updates has been enhancing numerical stability, preventing memory leaks, and generalizing the framework to support additional datasets. While the system is notably more robust, a few experimental features and auxiliary documentation files have been removed.

\section*{Key Improvements}
\subsection*{1. Multi-Dataset Support (CHB-MIT Expansion)}
\begin{itemize}[label=\textcolor{blue}{$\bullet$}]
    \item \textbf{Dataset Generalization:} The codebase now fully supports the \textbf{CHB-MIT} dataset alongside the original TUSZ dataset. 
    \item \textbf{Configuration Updates:} Added \texttt{CHBMIT\_INCLUDED\_CHANNELS} inside \texttt{constants.py} and expanded dataset \texttt{args.py} choices.
    \item \textbf{Marker Generation:} The hardcoded data split scripts were replaced with a versatile \texttt{generate\_markers.py} script capable of handling variable time windows (10s, 12s, 60s) and patient-level cross-validation splits.
\end{itemize}

\subsection*{2. Numerical Stability and Crash Prevention}
\begin{itemize}[label=\textcolor{blue}{$\bullet$}]
    \item \textbf{NaN/Inf Handlers:} Massive improvements to input stability. Corrupt or noisy EEG segments that result in NaNs or Infs are now aggressively caught using \texttt{torch.nan\_to\_num} in \texttt{main.py} and \texttt{EvoBrain.py}, preventing catastrophic training crashes.
    \item \textbf{Laplacian Eigenvector Fix:} The previously unstable Positional Encoding (PE) Laplacian decomposition was refactored. It now features safe fallbacks: attempting \texttt{get\_laplacian}, falling back to SVD on failure, and defaulting to zero-padding.
    \item \textbf{Loss Function Robustness:} Replaced a hardcoded assumption of BCE targets with dynamic type casting (\texttt{y.long()}) to properly support \texttt{CrossEntropyLoss}.
\end{itemize}

\subsection*{3. Memory Optimization}
\begin{itemize}[label=\textcolor{blue}{$\bullet$}]
    \item \textbf{Hidden State Leak Fixed:} In the original repository, \texttt{hidden\_all} states were being collected during both training and validation loops, leading to severe Out-Of-Memory (OOM) errors over long epochs. This was fixed by restricting hidden state logging explicitly to test sets (\texttt{is\_test and eval\_set == 'test'}).
\end{itemize}

\subsection*{4. Metrics and Hygiene Tools}
\begin{itemize}[label=\textcolor{blue}{$\bullet$}]
    \item \textbf{New Metrics:} Added tracking for Balanced Accuracy, Specificity, and PR-AUC. Specificity calculation was updated to avoid divide-by-zero errors.
    \item \textbf{Data Pipeline Tools:} Added \texttt{clean\_nan\_files.py} to actively quarantine corrupt \texttt{.pkl} files and \texttt{diagnose\_channels.py} to troubleshoot missing channels in raw EDF files.
\end{itemize}

\subsection*{5. Documentation Rewrite}
\begin{itemize}[label=\textcolor{blue}{$\bullet$}]
    \item The main \texttt{README.md} was completely overhauled. It now provides a clean, step-by-step replication pipeline covering raw EDF resampling, dataset splitting, training, and testing. It also includes hardware recommendations for memory constraints.
\end{itemize}

\section*{Downgrades and Removed Features}
\subsection*{1. Removed Architectures}
\begin{itemize}[label=\textcolor{red}{$\bullet$}]
    \item \textbf{Dense Inception Dropped:} The experimental \texttt{model/dense\_inception} directory (which included \texttt{dense\_inception.py} and \texttt{inceptions.py}) was removed entirely. If this was used as a baseline model, it is no longer available.
\end{itemize}

\subsection*{2. Lost Documentation}
\begin{itemize}[label=\textcolor{red}{$\bullet$}]
    \item \textbf{Specific Guides Removed:} The \texttt{model\_explanations.md} (detailing theoretical architectures) and \texttt{RTX4050\_RUN\_GUIDE.md} (hardware specific run configs) were deleted. While the main README improved, deep-dive context was lost.
\end{itemize}

\subsection*{3. Minor Regressions}
\begin{itemize}[label=\textcolor{orange}{$\bullet$}]
    \item \textbf{Argparse Typo:} In \texttt{args.py}, the default value for \texttt{--max\_seq\_len} was changed from an integer (\texttt{12}) to a string (\texttt{'12'}). While argparse often circumvents this during parsing, it is poor practice and can cause type errors if the attribute is accessed programmatically before CLI parsing.
\end{itemize}

\section*{Conclusion}
The transition from \texttt{EvoBrain-main} to \texttt{EvoBrain} represents a necessary shift from an experimental research codebase to a much more stable, multi-dataset pipeline. The memory fixes and numerical stability patches are highly commendable, though users porting over from the old repository should be aware of the dropped baseline architectures and missing theoretical documentation.

\end{document}
